\documentclass[12pt, a4paper]{article}
\usepackage[utf8]{inputenc}
\usepackage[T1]{fontenc}
\usepackage[margin=1in]{geometry}
\usepackage{tcolorbox}
\usepackage{setspace}
\usepackage{microtype}

\begin{document}

\begin{center}
    \LARGE\textbf{Four-Week Progress Report (SRFP 2026)}
\end{center}
\vspace{0.2cm}


\begin{tcolorbox}[
    colback=black!3,       
    colframe=black!70,     
    boxrule=0.8pt,         
    arc=4pt,              
    top=12pt,
    bottom=12pt,
    left=15pt,
    right=15pt
]
    \large
    \textbf{Student:} Devang Sachin Gupta (ENGS1253) \\
    \textbf{Institute:} Vellore Institute of Technology, Andhra Pradesh \\
    \textbf{Host Institute:} IIT Hyderabad \\
    \textbf{Supervisor:} Dr. G. V. V. Sharma
\end{tcolorbox}

\vspace{0.4cm}

\setstretch{1.3}

The first four weeks of my Summer Research Fellowship Programme (SRFP) internship at IIT Hyderabad under the supervision of Dr. G. V. V. Sharma have provided me with a strong foundation in the tools, concepts, and hardware platforms required for the research work to be undertaken during this internship. The activities carried out during this period were focused on developing my skills in Linux-based operating systems, digital electronics, embedded systems, and technical documentation.

I completed training in LaTeX{} and gained practical experience in preparing technical documents and reports in a professional academic format. Alongside this, I worked extensively with Linux systems, including the installation and configuration of a dual-boot environment. This experience enhanced my understanding of system administration, command-line utilities, package management, and the Linux development workflow. I also became familiar with troubleshooting common system-level issues and configuring development environments for embedded applications.

To improve my productivity in terminal-based environments, I learned and regularly used tools such as Termux, Neovim, and Ranger. Through various exercises and assignments, I developed proficiency in file management, text editing, Linux navigation, and command-line operations. Additionally, I explored looping constructs and shell-based automation techniques, which are valuable for efficient software development and research workflows. These activities helped me become more comfortable with a keyboard-driven development environment and improved my overall efficiency while working on technical tasks.

On the hardware side, I performed experiments involving the 7447 BCD-to-seven-segment decoder and the 7474 D flip-flop integrated circuit. These activities strengthened my understanding of combinational and sequential logic design. In parallel, I solved several GATE-level digital electronics problems related to these topics, reinforcing the theoretical concepts behind the hardware implementations and improving my analytical and problem-solving skills. The combination of theoretical study and practical experimentation provided a deeper understanding of digital system design.

I also gained hands-on experience in soldering and electronic prototyping. Furthermore, I initiated work on a digital clock project, through which I learned the fundamentals of hardware assembly, enclosure design, and 3D printing. This project has provided practical exposure to the complete process of developing and integrating hardware components into a functional system.

Overall, the first month of the internship has provided a comprehensive introduction to the tools, techniques, and hardware systems relevant to my research activities. The knowledge and practical experience gained during this period will serve as a strong foundation for the advanced hardware development, FPGA-based experimentation, and project work planned for the remaining duration of the internship.

\end{document}
